\begin{textsample}{POS Dim 4 – pa – Score 17.00 – pa/pa\_em\_95.txt}  \label{ex:f4_pos_004}
Um dos aspectos relacionados a esse declínio se deve a novos projetos, modalidades e/ou sistemas de ensino que passaram a ofertar o ensino médio.
Outros fatores são as dificuldades encontradas na infraestrutura ofertada pela escola de funcionamento ( municipal ) e o fornecimento de transporte regular que facilita o deslocamento desse aluno para a sede do município, lugar com maior oportunidade de emprego e lazer.
Contudo, o SOME continua com um expressivo atendimento nas localidades dos interiores do \textbf{estado}, sendo o sistema de ensino \textbf{predominante} na oferta do ensino médio, caracterizando -se “ (... ) como de suma importância para a educação, por ser um processo de democratização de ensino, onde oportuniza que muitos jovens almejem o ensino superior (... ) ” ( GAIA ; MENDES, 2020, p. 249 ).
Como se nota, nos escores apresentados acima, o SOME : > (... ) busca atender um público que naturalmente ficaria sem o acesso ao Ensino Médio, pela especificidade geográfica do \textbf{estado} paraense em que muitas vilas e comunidades rurais e \textbf{ribeirinhas} são de difícil acesso, além das condições financeiras das famílias em não poderem enviar seus filhos para estudar nos centros urbanos ( PEREIRA apud GAIA ; MENDES, 2020, p. 251 ) Com relação aos índices de rendimento escolar dos alunos do SOME no período em questão, constata -se a melhoria desses indicadores no gráfico 3.
A expansão do quantitativo de alunos que foram aprovados nas três séries do ensino Médio de 72,33 \% em 2014 para 75,04 \% em 2015, isto é, um crescimento 71 \% na taxa de aprovação escolar.
No que tange ao índice de reprovação não houve mudança significativa, permanecendo no patamar médio de 10 \%, a exceção do ano letivo de 2018, que ficou com o percentual de 4,97 \%.
Sobre o abandono escolar, percebe -se uma redução mínima de 0,98 \% entre os anos letivos de 2018 e 2019.
Por meio da apresentação dos dados de APROVAÇÃO e ABANDONO do ano letivo de 2020 exposta no gráfico 3, nota -se que não houve nenhum percentual de REPROVAÇÃO, tal constatação se deve ao fato do acometimento da pandemia mundial causada pelo vírus COVID-19, que atingiu também o \textbf{estado} do Pará.
Esta situação sanitária impulsionou a criação da Resolução de n{\EmojiFont °} 20 de 18 de janeiro de 2021, do Conselho Estadual de Educação do Pará que, em seu artigo 2o, §2o determina que : > Poderão ser aprovados os estudantes concluintes dos Ensinos Fundamental e Médio no ano letivo de 2020 que tiverem integralizado 75 \% da carga horária da respectiva série/ano da etapa da Educação Básica, sem prejuízo do alcance das competências e objetivos de aprendizagem relacionados à BNCC, garantindo -se a possibilidade de mudança de nível/etapa e de acesso ao Ensino Médio, Cursos Técnicos ou à Educação Superior, conforme o caso. ( CEE, 2021 ) Ressaltamos que a determinação do Conselho foi estendida a todos os alunos da educação básica da Secretaria de Estado de Educação.
Segundo o Sistema de Informação e Gestão do Pará, no ano de 2020 não tivemos reprovados e a taxa de abandono foi quase nula.
Alguns fatores podem ter contribuído para essa diferença em relação aos anos anteriores, tal como a Portaria do Conselho elencada acima e a política de distribuição do cartão de vale alimentação realizada pelo Governo do \textbf{Estado} do Pará, que beneficiou todos os alunos matriculados na rede estadual de ensino.
Para uma reestruturação curricular há que se considerar uma mudança no SOME que passe a buscar uma nova postura docente, pautada pelo pensamento autônomo e \textbf{libertador} do educando, neste sentido : > A educação que se impõe aos que verdadeiramente se comprometem com a libertação não pode fundar -se numa compreensão dos \textbf{homens} como seres “ vazios ” a quem o mundo “ encha ” de conteúdos ; não pode basear -se numa consciência espacializada, mecanicamente compartimentada, mas nos \textbf{homens} como “ corpos conscientes ” e na consciência como consciência intencionada ao mundo.
Não pode ser a do depósito de conteúdos, mas a da problematização dos \textbf{homens} em suas relações com o mundo ( FREIRE, 1987, p. 67 ).
Nesse contexto, a educação precisa ser dialógica e fomentadora da práxis, que necessariamente “ [... ] implica ação e a reflexão dos \textbf{homens} sobre o mundo para transformá -lo ” ( FREIRE, 1987, p. 67 ).
Assim, o SOME envidará esforços no sentido de promover uma educação problematizadora nas comunidades às quais atende, tonalizada pelas realidades de cada lugar, bem como pela autenticidade do pensar de educador e educandos, opondo -se, sobremaneira, à “ educação bancária ”, que segundo Freire ( 1987 ) trata de uma concepção mecânica da consciência, considerada como um repositório vazio, que precisa ser preenchido, restando aos alunos aceitar passivamente o depósito de conteúdo, encarados por meio de uma única visão, que objetiva apassivar ainda mais os educandos.
Freire nos provoca quanto ao uso da educação como lugar de reflexão para intervenção no mundo, por meio do viés histórico como viabilidade, por isso : > Só na história como possibilidade e não como determinação se percebe e se vive a subjetividade em sua dialética relação com o objetividade.
É percebendo e vivendo a história como possibilidade que experimento plenamente a capacidade de comparar, de ajuizar, de escolher, de decidir, de romper.
E é assim que \textbf{mulheres} e \textbf{homens} eticizam o mundo, podendo, por outro lado tornar -se transgressores da própria ética. ( FREIRE, 2000, pág. 57 ) Nessa direção, à educação não é facultada a decisão de ser \textbf{neutra}, pois pode servir tanto aos princípios \textbf{libertadores}, de inserção crítica do ser humano no mundo, de rupturas, tornando -se consciente da presença do outro e de si no mundo, quanto a aceitação passiva de estruturas historicamente injustas e desiguais.
Em busca de uma reestruturação para além da matriz curricular, o SOME tem procurado (re)conhecer os sujeitos coletivos \textbf{populares}, bem como suas experiências, produção de conhecimentos, de culturas e histórias como significativas, compreendendo suas relevâncias como sujeitos da história, dando vez e voz à estes coletivos \textbf{populares} na relação pedagógica.
Tendo em vista que há um perigo danoso em se considerar uma única história, em observar de maneira unilateral os coletivos humanos, em razão de nossa tradição política e cultural que com base segregadora instituiu uma polaridade de coletivos, evidenciando uns em detrimento de outros, determinando com isso ausências de sujeitos sociais no \textbf{território} do conhecimento.
A este respeito Arroyo ( 2011, p. 139 ) sugere que : > Um caminho para avançar pode ser entender melhor os processos históricos que vêm de longe e se perpetuam de não reconhecimento de coletivos \textbf{populares}, dos trabalhadores como produtores de saberes, culturas, modos de pensar.
Como sujeitos de história.
Questões que tocam em cheio a docência e os currículos.
O SOME tem em seu perfil um diversificado \textbf{território} de multiplicidades, posto que atua com alunos oriundos de \textbf{quilombos}, assentamentos, produtores rurais, \textbf{indígenas}, \textbf{ribeirinhos}, extrativistas, pescadores, entre outros e é em razão dessa vasta diversificação identitária que não se deixa aprisionar, por isso necessita de um processo educativo que considere a centralidade dos seres humanos como sujeitos da história, pois “ todos os conhecimentos sustentam práticas e constituem sujeitos ” ( SANTOS apud ARROYO, 2011, p. 139 ).
O ano de 2020 registrou 34.580 alunos matriculados no Sistema Modular por meio do Sistema de Informação de Gestão Escolar do Pará ( SIGEP ), que atende massivamente ao alunado que tonaliza o ensino médio nas comunidades do interior do \textbf{Estado}.
Este perfil do atendimento no SOME exige que cada vez mais se busque o reconhecimento da diversidade dos sujeitos jovens, que têm outros modos de pensar e diversas formas de ver o mundo.
Nesse sentido, valem as observações relativas à necessidade de uma nova organização curricular, acompanhada de programas de formação inicial e continuada de professores, mais adequada à nova realidade do jovem que vive em relações sociais e produtivas marcadas pela \textbf{exclusão}, pela ausência de projeto de futuro, pela complexidade tecnológica e dos meios de comunicação, pela flexibilidade, pela instabilidade, pela intensificação e pelo estresse, para citar apenas as dimensões mais evidentes ( KUENZER, 2010, p. 869 ).
Dentro da perspectiva de Freire ( 1987 ) e da concepção curricular do Novo Ensino Médio do Pará, o público alvo do SOME, em sua maioria na faixa etária de 15 à 18 anos de idade, contará com uma proposta curricular integral, tendo como aliada a interdisciplinaridade na construção do conhecimento globalizado e como articuladora na contextualização dos saberes, considerando as múltiplas realidades desse público.
O ambiente educacional incentivador do protagonismo juvenil é permeado de possibilidades de aprendizagem que favorecem a contextualização da realidade local e global num ambiente reflexivo e participativo para manifestações, debates, pesquisas, construções individuais e coletivas no âmbito escolar e social.
Problematizar a realidade e as ações necessárias para a construção de um projeto de vida consciente e \textbf{libertador}, no qual a consciência de sujeito da própria vida é importante para motivar esse jovem a ser um agente transformador de sua realidade.
Outra particularidade do SOME é a necessária imersão e adequação dos conhecimentos às realidades das comunidades amazônidas, alinhando -se aos princípios norteadores da educação básica paraense quando ao eleger o respeito às diversas culturas \textbf{amazônicas} e suas inter-relações no espaço e no tempo como princípio.
Este princípio contribui para a centralidade nos currículos da produção histórica e cultural de \textbf{homens} e \textbf{mulheres} da Amazônia, refletidas sobre patrimônio material, imaterial, ambiental, nas danças, nas festividades \textbf{populares} e \textbf{religiosas}, nos costumes, no \textbf{artesanato}, na produção artística e literária, na culinária, na produção \textbf{agrícola} e na riqueza mineral, proporcionando uma pluridiversidade e uma rica interculturalidade.
Nesse caminhar, os princípios do documento curricular do \textbf{Estado} do Pará ( 2019 ) que tratam do respeito às diversas culturas \textbf{amazônicas} e suas inter-relações no espaço e no tempo, educação para a sustentabilidade ambiental, social e econômica, bem como a interdisciplinaridade no processo ensino-aprendizagem devem compor a arquitetura curricular do SOME.
Abordar tamanha riqueza de trajetórias identitárias é um desafio a ser superado coletiva e cotidianamente no modular, sobretudo porque historicamente diversos interesses e grupos constituíram e constituem esses sujeitos, por meio de relações de poder que estão imbricadas em todas as instâncias sociais ( LOURO, 1997 ), porém, não podemos deixar de levar em consideração que estes sujeitos singulares conformam ou não tais preceitos.
Portanto, existem grupos que a partir das relações de poder, determinam como “ o outro ” vai ser representado e quem vai ser objeto dessa representação.
Em face dessas representações sociais, imersas em relações de poder, os sujeitos também vão se construindo.
Neste \textbf{bojo}, Louro ( 1997, p. 102, grifos da autora ) nos alerta para “ [... ] o fato de que os significados que as representações acabam produzindo não preexistem no mundo, mas eles têm que ser criados, e são criados socialmente, isto é, são criados através de relações sociais de poder ”.
O poder para Michel Foucault ( 1979 ) aparece como um mecanismo com dois grandes lados : um que institui interdição, sujeição, \textbf{exclusão} e dominação e outro que provoca produção, incitação e reação.
A propósito disso, Neto ( 2004. p. 04 ) aborda que não existe um poder com características indestrutíveis, que nos tome por completo, eliminando a capacidade humana de qualquer probabilidade de insurreição, posto que segundo a autora “ (... ) a capacidade de se insurgir, de se rebelar, de resistir são pontos construtivos das relações de poder, pois estão presentes e complementam o campo de forças onde as relações de poder e os pontos de resistência se entrelaçam ”.
Como bem afirma Foucault ( 1988 ) relações de poder e possibilidades de resistência não estão situados em lugares específicos, não são privilégios deste ou daquele grupo e nem podem ser encarados como indiferentes aos acontecimentos sociais.
Mas, sobretudo, estão enredados com a (re)produção de saberes, de mudanças e transformações. > [... ] lá onde há poder há resistência e, no entanto ( ou melhor, por isso mesmo ) esta nunca se encontra em posição de exterioridade em relação ao poder. [... ] Não existe, com respeito ao poder um lugar da grande Recusa-alma da revolta, foco de todas as rebeliões, leitura do revolucionário.
Mas sim, resistências, no plural. [... ] As resistências não se reduzem a uns poucos princípios heterogêneos ; mas não é por isso que seja ilusão, ou promessa necessariamente desrespeitada.
Elas são o outro termo nas relações de poder, inscrevem -se nessas relações como interlocutor irredutível ( FOUCAULT, 1998, p. 91-92 ).
Sendo assim, poder e resistência são os inseparáveis lados de uma mesma moeda.
São fios que se cruzam na tecitura de um tecido, consolidando pontos firmes que ajudam na composição do corpo social em toda sua extensão.
São enredos que nos compõem, marcam nossos espaços, constituem nossos gestos e ditam nossos comportamentos.
Na transposição da perspectiva teórico-metodológica foucaultiana para o campo da discussão em torno da formação e currículo podemos inferir que professoras/es e alunas/os não estão apenas sujeitos a uma formação ou a um currículo que determinam, peremptoriamente, seu fazer, mas ao contrário, estes sujeitos experienciam as relações e práticas de poder em seus espaços de atuação sejam eles profissionais ou pessoais.
Essas pessoas (re)criam discursos, práticas, conhecimentos, espraiam suas posições políticas, seus preconceitos, suas pré-noções, descobrem táticas de reivindicações e sobrevivência, mesmo quando são expostos a condições inadequadas de trabalho, salário, moradia, estudo, entre outros.
Ao fazer um balanço geral das relações de poder e resistência, o grande desafio é superar as omissões históricas que se materializaram nas perspectivas curriculares aplicadas no ensino médio \textbf{amazônico}, especialmente no SOME.
Um currículo que considere as realidades locais implica refletir principalmente nas práticas e concepções históricas, no \textbf{território} e suas \textbf{territorialidades}, que considere a aprendizagem decorrente da realidade vivida e, consequentemente, do direito de repensar o sentido/significado de estar no mundo e de ressignificar os processos de ensino-aprendizagem capaz de transformar essa mesma realidade em interrelação com a interculturalidade.

% matched lemmas: agrícola, amazônico, artesanato, bojo, estado, exclusão, homem, indígena, libertador, mulher, neutro, popular, predominante, quilombo, religioso, ribeirinho, territorialidade, território
\end{textsample}
